\documentclass[12pt]{article}

\usepackage{sbc-template}
\usepackage{graphicx,url}
\usepackage[utf8]{inputenc}
\usepackage[brazil]{babel}
%\usepackage[latin1]{inputenc}  

     
\sloppy

\title{K-ORCH: ferramenta para controle e manejo de sistemas multikernel}

\author{Eduardo Kaluf\inst{1}, Carlos Alberto Maziero\inst{2}}

\address{
  Departamento de Informática -- Universidade Federal do Paraná (UFPR)\\
  Curitiba -- PR -- Brazil
  \email{eduardo@kaluf.com.br, maziero@inf.ufpr.br}
}

\begin{document} 

\maketitle

\begin{abstract}
  The sharp increase in the number of cores per processor, combined with the growing heterogeneity of the hardware in use,
  ends up exposing the trade-offs and limitations of traditional monolithic operating system architectures, visible in
  memory hierarchies, interconnects, and I/O configurations, as well as in the management of internal data and control structures.

  One line of research proposes, as an alternative, running multiple instances of a subset of kernels, with different
  configurations, on the same physical processor, all of them operating much like a distributed system.
  This approach to the OS is now known as the multikernel, and it has been shown to be viable as a solution that goes beyond theory.

  However, orchestrating distributed nodes that use the multikernel architecture remains a challenge, due to the inherent
  complexity of dealing with partitioned processors and to the relative novelty of the architecture.
  This work proposes the development of an orchestrator specifically designed for multikernel environments, making it possible
  to extend the architecture beyond a single local machine, with dispatching, load balancing, and monitoring of the active instances in the environment.
\end{abstract}

\begin{resumo}
  O grande aumento da quantidade de núcleos por processador, em conjunto com a crescente heterogeneidade do hardware utilizado,
  acaba por expor contrapartidas e limites das arquiteturas tradicionais de sistemas operacionais monolíticos, visíveis nas
  hierarquias de memória, interconexões e configurações de E/S, além do gerenciamento das estruturas de dados e de controle internas.

  Uma linha de pesquisa apresenta como alternativa executar múltiplas instâncias de um subconjunto de kernels, com diferentes
  configurações, em um mesmo processador físico, onde todos funcionam de forma similar a um sistema distribuído.
  Essa maneira de tratar o SO hoje é conhecida como multikernel e já se demonstrou viável como uma solução que vai além apenas da teoria.

  Porém, a orquestração de nós distribuídos que utilizam a arquitetura multikernel ainda é um desafio, pela complexidade inerente
  de lidar com processadores particionados e pela natureza recente da arquitetura.
  Este trabalho propõe o desenvolvimento de um orquestrador específico para ambientes multikernel, dando a possibilidade
  de evoluir a arquitetura para além de uma máquina local, com despacho, balanceamento e monitoramento das instâncias ativas no ambiente.
\end{resumo}


%\section{General Information}
%
%All full papers and posters (short papers) submitted to some SBC conference,
%including any supporting documents, should be written in English or in
%Portuguese. The format paper should be A4 with single column, 3.5 cm for upper
%margin, 2.5 cm for bottom margin and 3.0 cm for lateral margins, without
%headers or footers. The main font must be Times, 12 point nominal size, with 6
%points of space before each paragraph. Page numbers must be suppressed.
%
%Full papers must respect the page limits defined by the conference.
%Conferences that publish just abstracts ask for \textbf{one}-page texts.
%
%\section{First Page} \label{sec:firstpage}
%
%The first page must display the paper title, the name and address of the
%authors, the abstract in English and ``resumo'' in Portuguese (``resumos'' are
%required only for papers written in Portuguese). The title must be centered
%over the whole page, in 16 point boldface font and with 12 points of space
%before itself. Author names must be centered in 12 point font, bold, all of
%them disposed in the same line, separated by commas and with 12 points of
%space after the title. Addresses must be centered in 12 point font, also with
%12 points of space after the authors' names. E-mail addresses should be
%written using font Courier New, 10 point nominal size, with 6 points of space
%before and 6 points of space after.
%
%The abstract and ``resumo'' (if is the case) must be in 12 point Times font,
%indented 0.8cm on both sides. The word \textbf{Abstract} and \textbf{Resumo},
%should be written in boldface and must precede the text.
%
%\section{CD-ROMs and Printed Proceedings}
%
%In some conferences, the papers are published on CD-ROM while only the
%abstract is published in the printed Proceedings. In this case, authors are
%invited to prepare two final versions of the paper. One, complete, to be
%published on the CD and the other, containing only the first page, with
%abstract and ``resumo'' (for papers in Portuguese).
%
%\section{Sections and Paragraphs}
%
%Section titles must be in boldface, 13pt, flush left. There should be an extra
%12 pt of space before each title. Section numbering is optional. The first
%paragraph of each section should not be indented, while the first lines of
%subsequent paragraphs should be indented by 1.27 cm.
%
%\subsection{Subsections}
%
%The subsection titles must be in boldface, 12pt, flush left.
%
%\section{Figures and Captions}\label{sec:figs}
%
%
%Figure and table captions should be centered if less than one line
%(Figure~\ref{fig:exampleFig1}), otherwise justified and indented by 0.8cm on
%both margins, as shown in Figure~\ref{fig:exampleFig2}. The caption font must
%be Helvetica, 10 point, boldface, with 6 points of space before and after each
%caption.
%
%%\begin{figure}[ht]
%%\centering
%%\includegraphics[width=.5\textwidth]{fig1.jpg}
%%\caption{A typical figure}
%%\label{fig:exampleFig1}
%%\end{figure}
%
%%\begin{figure}[ht]
%%\centering
%%\includegraphics[width=.3\textwidth]{fig2.jpg}
%%\caption{This figure is an example of a figure caption taking more than one
%%  line and justified considering margins mentioned in Section~\ref{sec:figs}.}
%%\label{fig:exampleFig2}
%%\end{figure}
%
%In tables, try to avoid the use of colored or shaded backgrounds, and avoid
%thick, doubled, or unnecessary framing lines. When reporting empirical data,
%do not use more decimal digits than warranted by their precision and
%reproducibility. Table caption must be placed before the table (see Table 1)
%and the font used must also be Helvetica, 10 point, boldface, with 6 points of
%space before and after each caption.
%
%%\begin{table}[ht]
%%\centering
%%\caption{Variables to be considered on the evaluation of interaction
%%  techniques}
%%\label{tab:exTable1}
%%\includegraphics[width=.7\textwidth]{table.jpg}
%%\end{table}
%
%\section{Images}
%
%All images and illustrations should be in black-and-white, or gray tones,
%excepting for the papers that will be electronically available (on CD-ROMs,
%internet, etc.). The image resolution on paper should be about 600 dpi for
%black-and-white images, and 150-300 dpi for grayscale images.  Do not include
%images with excessive resolution, as they may take hours to print, without any
%visible difference in the result.
%
%\section{References}
%
%Bibliographic references must be unambiguous and uniform.  We recommend giving
%the author names references in brackets, e.g. \cite{knuth:84},
%\cite{boulic:91}, and \cite{smith:99}.
%
%The references must be listed using 12 point font size, with 6 points of space
%before each reference. The first line of each reference should not be
%indented, while the subsequent should be indented by 0.5 cm.

\bibliographystyle{sbc}
\bibliography{sbc-template}

\end{document}
